\documentclass[12pt,a4paper]{article}

\usepackage[utf8]{inputenc}
\usepackage[T1]{fontenc}
\usepackage[spanish]{babel}
\usepackage{amsmath, amssymb}
\usepackage{geometry}
\usepackage{setspace}
\usepackage{parskip}
\usepackage{enumitem}
\usepackage{hyperref}

\geometry{margin=2.5cm}
\onehalfspacing

\title{\textbf{Hito 1 (draft)}}
\author{
    \textbf{Nombre Alumno:} Eduardo Caballero \\
    \textbf{Nombre Profesor guía:} Sebastián Cea
}
\date{}

\begin{document}

\maketitle

% ============================================================
\section*{Introducción}
% ============================================================

El presente documento revisa la tesis sobre \textit{Financial Networks \& Technological
Innovation: A Dual-Network Model} (Venegas, 2024), que aborda cómo las restricciones
regulatorias en redes financieras impiden relaciones bilaterales eficientes y generan
pérdidas de bienestar. La relevancia del problema es significativa: las infraestructuras
financieras alternativas (NBFI o \textit{shadow banking}) representan hoy USD~238 billones a
nivel global, equivalentes al 49,1\% de los activos financieros mundiales (Financial
Stability Board, 2024).

Con base en la lectura de \textit{``Financial Networks \& Technological Innovation: A Dual-Network
Model''} (Venegas, 2024) se destacarán los puntos base sobre los que la tesis se sostiene, los
elementos centrales del problema, el contexto en el que se enmarca y la extensión que realiza
la tesis sobre el modelo de \textit{``Strategic Negotiations in Endogenous Network Formation''}
(Jalan y Chakrabarti, 2024). Además, se expondrán algunas oportunidades de mejora al modelo.
El documento se dividirá en tres grandes capítulos siguiendo el orden establecido en la tesis:
Marco teórico, modelo y resultados.

La tesis estudia la alternativa de una infraestructura paralela a la estructura financiera que
opera con regulaciones diferentes a las instituciones tradicionales. Luego se introducen
distintos modelos estudiados y referenciados en el marco teórico donde se demarcan sus
limitaciones. Finalmente se introduce un modelo \textit{dual-network infrastructure} y se
demuestra que la coexistencia de ambos sistemas con diferentes regulaciones mejora el
bienestar y no elimina al sistema tradicional. El resultado central es que la sola competencia
entre intermediarios dentro de la red restringida recupera el bienestar parcialmente (deja un
gap de 27,9\% respecto del benchmark sin restricción), mientras que habilitar una
infraestructura alternativa logra una recuperación completa del bienestar. Se modelan a los
agentes en redes paralelas, donde cada uno adjudica sus propios contratos para maximizar el
bienestar:

\begin{itemize}
    \item $T$: Banca tradicional con sus debidas restricciones.
    \item $A$: Sistema alternativo que opera bajo distintos tratamientos.
\end{itemize}


% ============================================================
\section*{Marco Teórico}
% ============================================================

El capítulo presenta el marco teórico que sustenta el análisis de redes financieras y la
competencia entre infraestructuras tradicionales y alternativas. Se definen las redes como
conexión entre nodos y links de acuerdo con la siguiente forma: los agentes (nodos) toman
decisiones estratégicas sobre con quiénes formar links considerando los beneficios esperados
y sus costos. Con base en estas redes se pueden modelar fenómenos como el contagio
financiero, la propagación del riesgo sistémico y la estabilidad del sistema económico.

Se hace referencia a 3 modelos de redes financieras:

\begin{itemize}
    \item \textbf{Eisenberg y Noe (2001):} Redes financieras a través de matriz de obligaciones
    que representa de buena forma el efecto cascada en la red cuando ocurre un \textit{default}.

    \item \textbf{Elliott, Golub y Jackson (2014):} Modelo de red para contagio y cómo puede
    amplificar o mitigar riesgo sistemático.

    \item \textbf{Jalan y Chakrabarti (2024):} Introduce una matriz $Q$ simétrica y semidefinida
    positiva que captura la estructura de costos y beneficios entre pares de agentes, junto con
    una función de utilidad media-varianza que permite modelar decisiones óptimas de
    contratación bilateral en equilibrio \textit{pairwise stable}. Es el modelo base que la tesis
    extiende.
\end{itemize}

Sin embargo, las redes tradicionales presentan limitaciones relevantes, tales como altos costos
de intermediación, la opacidad en la información, el riesgo sistémico derivado de la
centralización y la baja inclusión financiera. Además, los efectos de red pueden generar tanto
beneficios, como mayor acceso a financiamiento, como costos, asociados a congestión y
fricciones operacionales.

Posteriormente, el capítulo analiza la competencia entre infraestructuras financieras desde la
perspectiva de mercados de dos lados y teoría de plataformas. De esta forma, las redes
funcionan como plataformas que conectan a tanto oferentes y demandantes de liquidez, donde
los efectos de red y las decisiones de \textit{multi-homing} determinan el equilibrio competitivo.
Las referencias en la literatura muestran que la coexistencia de múltiples plataformas depende
de factores tales como: costos de cambio, la diferenciación entre servicios y la presencia de
externalidades de red.

Un elemento central es el concepto de arbitraje regulatorio, donde las restricciones impuestas
a las redes tradicionales incentivan la aparición de otras infraestructuras alternativas que
permiten evadir las limitaciones impuestas. Por ejemplo, se ha demostrado que una parte
importante del crecimiento del Fintech y del \textit{shadow banking} responde justamente a
estas restricciones más que a sus ventajas tecnológicas intrínsecas.

Asimismo, el capítulo conecta el análisis de redes con el equilibrio general y la provisión de
liquidez. Para lograr el equilibrio general, las redes financieras determinan cómo se asigna el
riesgo entre los distintos agentes y cómo los precios de los activos reflejan la estructura de
las relaciones bilaterales. Respecto a la liquidez, las redes operan como canales que la proveen
y la redistribuyen entre los agentes. Esto se traduce en que en redes más densas la liquidez
fluye con menor fricción, pero también se amplifica la transmisión de \textit{shocks}.

El capítulo también aborda la adopción tecnológica en presencia de efectos de red; aquí se
destacan problemas de coordinación y posibles situaciones de \textit{lock-in}, donde
tecnologías inferiores pueden persistir debido a su ventaja histórica y a que el mercado ya las
adoptó. Asimismo, se analizan los costos de cambio, que incluyen costos de transacción,
aprendizaje y coordinación, los cuales influyen en la velocidad y magnitud de adopción de
nuevas infraestructuras.

También se discuten las condiciones bajo las cuales pueden coexistir múltiples redes o surgir
dinámicas de tipo \textit{winner-takes-all}. La evidencia teórica sugiere que la coexistencia
es posible cuando existen diferencias regulatorias, beneficios y costos balanceados,
heterogeneidad entre agentes y complementariedades entre plataformas. En contraste, la
dominancia de una sola red ocurre cuando una infraestructura presenta ventajas claras o
cuando los costos de adopción de alternativas son prohibitivos.

Por su parte, la tecnología \textit{blockchain} como base de las finanzas descentralizadas, se
presenta como una infraestructura alternativa, ya que permite realizar transacciones sin
intermediarios mediante un registro distribuido y transparente. Sus principales beneficios son
la reducción de asimetrías de información, menores costos de transacción y mayor inclusión
financiera, además de la automatización de contratos a través de \textit{smart contracts}. No
obstante, su adopción enfrenta limitaciones relevantes, como costos operacionales asociados
a tarifas y congestión, problemas de escalabilidad, riesgos de privacidad y barreras
regulatorias y de adopción tecnológica que dificultan su implementación en el mercado masivo.


% ============================================================
\section*{Modelo}
% ============================================================

\subsection*{1. Modelo base de Jalan y Chakrabarti (2024)}

El modelo base que la tesis extiende considera una red de $N$ agentes que negocian contratos
bilaterales bajo una función de utilidad media-varianza. Cada agente $i$ posee un vector de
creencias $\mu_i$ sobre los retornos esperados de los contratos, una matriz de covarianzas
$\Sigma_i$ que captura el riesgo de su portafolio, y un coeficiente de aversión al riesgo
$\gamma_i$. La asignación que cada agente realiza a sus contratos se denota como un vector
$w_i$, y los pagos asociados se recogen en una matriz antisimétrica $P$ (lo que $i$ paga a $j$
es lo opuesto a lo que $j$ paga a $i$).

La utilidad de un agente $i$ toma la forma:
\begin{equation}
    g_i(W, P) = w_i^\top (\mu_i - P e_i) - \gamma_i\, w_i^\top \Sigma_i\, w_i,
\end{equation}
donde el primer término representa el retorno esperado neto de pagos y el segundo el costo
cuadrático asociado al riesgo. Cada agente elige $w_i$ para maximizar esta utilidad sujeta a
una restricción regulatoria capturada por una matriz de proyección $\Psi_i$, que determina qué
contratos puede o no puede tomar. La solución óptima admite la forma cerrada:
\begin{equation}
    w_i^* = Q_i (\mu_i - P e_i),
\end{equation}
donde la matriz $Q_i$ agrega aversión al riesgo, covarianza y restricciones regulatorias en un
único objeto.

\paragraph{Ejemplo 1: contrato de préstamo (Punto 3.1 en tesis).}
Considerando un contrato de préstamo entre dos agentes donde el prestatario $i$ recibe el
principal hoy y promete devolver capital más intereses en el futuro, su retorno esperado
$\mu_{ij}$ asociado al contrato podría ser, por ejemplo, 7\%. El prestamista $j$ entrega el
principal hoy y espera recibir el repago, por lo que su retorno esperado $\mu_{ji}$ desde su
perspectiva podría ser 3\%. Es importante notar que las creencias de $i$ y $j$ sobre el mismo
contrato no necesitan sumar cero: cada agente evalúa el contrato desde su propio punto de
vista, considerando su costo de oportunidad y la correlación con el resto de su portafolio. Esta
heterogeneidad es lo que abre espacio para que ambos agentes encuentren beneficio mutuo en
un contrato y negocien un pago $P_{ij}$ que reparta el excedente.

\paragraph{Ejemplo 4: pagos de incentivo entre dos firmas (Punto 3.3 en tesis).}
Un segundo ejemplo permite ilustrar por qué los pagos bilaterales son esenciales para alcanzar
un equilibrio. En este ejemplo se consideran dos firmas con distintas esperanzas sobre el
retorno conjunto, capturadas en una matriz:
\begin{equation}
    M = \begin{pmatrix} 0 & 3 \\ 1 & 4 \end{pmatrix},
\end{equation}
con matrices de covarianza $\Sigma_1 = \Sigma_2 = \mathrm{diag}(1, 2)$ y aversión al riesgo
$\gamma = 1$.

Sin pagos entre firmas, cada una resuelve su problema por separado: la firma~1 quiere asignar
$x^* = 1/4$ al contrato bilateral, mientras que la firma~2 quiere asignar $x^* = 3/2$. Como ambos
valores deben coincidir para que el contrato exista, no hay acuerdo posible y la red queda vacía.
Al permitir un pago $P$ entre firmas, las funciones objetivo se desplazan. Con $P = 5/3$, ambas
firmas convergen al mismo óptimo $x^* = 2/3$, y el contrato bilateral se vuelve estable. La
matriz $\mu_1$ pasa de $\begin{pmatrix}0\\1\end{pmatrix}$ a $\begin{pmatrix}0\\1+5/3\end{pmatrix}$.
Este ejemplo muestra que el rol de los pagos no es solo redistributivo, sino que es lo que
permite que firmas con creencias heterogéneas alineen sus asignaciones óptimas y que por lo
tanto exista un equilibrio bilateral.

\subsubsection*{Estabilidad y negociación \textit{pairwise}}

El modelo define una asignación $(W, P)$ como factible si $W$ es simétrica (lo que $i$
adjudica al contrato con $j$ coincide con lo que $j$ adjudica al mismo contrato), $P$ es
antisimétrica y ambos respetan las restricciones regulatorias $\Psi_i$. Una asignación es
estable si ningún par de agentes puede mejorar simultáneamente desviándose
unilateralmente, lo que corresponde a un equilibrio de Nash en negociaciones \textit{pairwise}.
Jalan y Chakrabarti (2024) muestran (Teorema 3.3.1) que un punto estable existe si y solo si
una condición de proyección sobre el espacio columna de las restricciones se satisface, y
proponen un algoritmo iterativo de negociaciones \textit{pairwise} que converge al equilibrio
bajo un parámetro de amortiguamiento $\eta$.

\subsection*{2. Restricción regulatoria y competencia dentro de la red}

El modelo base permite incorporar de forma natural restricciones regulatorias sobre los
contratos que cada agente puede tomar, a través de la matriz de proyección $\Psi_i$. Cuando
$\Psi_i$ prohíbe el contrato directo entre dos agentes, la única forma en que ambos pueden
vincularse es mediante un intermediario habilitado en ambos extremos. Esta estructura genera
rentas de intermediación que dependen del poder de mercado del intermediario y del grado de
competencia que enfrenta dentro de la red.

\paragraph{Ejemplo 5: red de tres firmas con prohibición regulatoria (Punto 3.3.4 en tesis).}
El siguiente ejemplo considera 3 agentes: un banco nacional ($N$), un banco local ($L$) y una
firma local ($F$). La regulación prohíbe el contrato directo entre el banco nacional y la firma
local (esto es, $\Psi$ desactiva la posición $(N, F)$ en la matriz de contratos), por lo que el
único canal disponible para que el banco nacional financie a la firma local es a través del banco
local. Como $L$ es el único intermediario habilitado, captura una renta monopólica: en el
equilibrio del modelo la firma local termina pagando aproximadamente $0{,}23$ al banco local,
mientras que el banco nacional paga aproximadamente $0{,}26$ al banco local por intermediar
la operación. La cadena de intermediación $N \to L \to F$ transfiere recursos, pero $L$ se queda
con el diferencial.

El equilibrio \textit{pairwise stable} converge a una asignación $W_\infty$ con aproximadamente
$[0{,}26,\, 0{,}56]$, donde el banco local concentra una posición desproporcionada relativa a
su rol funcional. Este ejemplo demuestra cómo una restricción regulatoria, en ausencia de
competencia, no solo impide la relación bilateral más eficiente, sino que además genera rentas
para el intermediario habilitado, deteriorando el bienestar agregado.

\paragraph{Extensión: incorporar competencia mediante una cuarta firma.}
Para aislar el efecto de la competencia entre intermediarios, la tesis introduce una cuarta firma
—que puede ser un intermediario de tipo Fintech— con un costo y un perfil de riesgo similares
al banco local. En este escenario el banco nacional puede ahora intermediar a través de $L$ o
de la Fintech, y ambos intermediarios compiten por el flujo.

El equilibrio resultante reparte el flujo aproximadamente en partes iguales entre los dos
intermediarios, lo que destruye la renta monopólica que $L$ capturaba en el caso de tres
firmas. Sin embargo, la ineficiencia estructural persiste ya que la firma local sigue sin poder
contratar directamente con el banco nacional porque la restricción regulatoria ($\Psi$ desactiva
$(N, F)$) no se ha removido, sino que se ha duplicado el número de intermediarios. La
competencia entre intermediarios elimina las rentas, pero no recupera el bienestar perdido por
la prohibición original. Este resultado motiva el salto al modelo de redes duales: si la
competencia dentro de la red restringida no basta, la pregunta natural es qué ocurre cuando se
permite competencia entre infraestructuras con regulaciones distintas.

\subsection*{3. Modelo \textit{dual-network}}

Cada agente $i$ conserva sus primitivos del modelo base pero ahora enfrenta dos matrices de
proyección regulatoria: $\Psi^T_i$ para la red tradicional y $\Psi^A_i$ para la red alternativa. La
asignación total que cada agente realiza al contrato con $j$ se descompone como:
\begin{equation}
    w_i^* = w^T_i + w^A_i,
\end{equation}
donde $w^T_i$ es la posición tomada en la red tradicional y $w^A_i$ la posición tomada en la red
alternativa. Cuando $\Psi^T_i \neq \Psi^A_i$ existe arbitraje regulatorio: contratos imposibles en
$T$ pueden ejecutarse en $A$.

La utilidad incorpora dos elementos nuevos asociados al uso de la red alternativa. Primero, un
beneficio de transparencia $\tau \geq 0$ que captura las ganancias de utilizar una
infraestructura como \textit{blockchain} (registro distribuido, menores asimetrías de
información, automatización vía \textit{smart contracts}). Segundo, un costo cuadrático de
fricción $\lambda_A \geq 0$ sobre la posición $w^A_i$, que captura los costos operacionales, de
adopción y de escalabilidad propios de la red alternativa. La utilidad extendida toma la forma:
\begin{equation}
    g_i(W, P) = w_i^\top (\mu_i + \tau e^A_i - P e_i) - \gamma_i\, w_i^\top \Sigma_i\, w_i
    - \lambda_A \|w^A_i\|^2,
\end{equation}
donde el término $\tau e^A_i$ premia el uso de $A$ y el término $\lambda_A \|w^A_i\|^2$ lo
penaliza, generando un \textit{trade-off} interno entre transparencia y fricción.

\subsubsection*{Existencia y consistencia con el modelo base}

El Teorema 3.5.6 extiende el resultado de existencia del modelo base al \textit{setting} dual:
bajo condiciones análogas sobre las restricciones combinadas $(\Psi^T, \Psi^A)$, existe un
punto estable en el equilibrio \textit{dual-network} y el algoritmo iterativo de negociaciones
\textit{pairwise} converge a él. La Proposición 3.5.7 establece la consistencia con el modelo
original. El modelo dual se reduce al modelo base de Jalan y Chakrabarti en tres casos límite:
\begin{itemize}
    \item Cuando $\lambda_A \to 0$ y $\tau = 0$ (la red alternativa no aporta ni cobra).
    \item Cuando $\Psi^A$ desactiva todos los contratos (la red alternativa no existe
    operativamente).
    \item Cuando $\Psi^T = \Psi^A$ (ambas redes son regulatoriamente idénticas y no hay
    arbitraje).
\end{itemize}
En estos casos $w^A_i = 0$ en el equilibrio y la asignación coincide con la del modelo de una
sola red.


% ============================================================
\section*{Resultados}
% ============================================================

El capítulo de resultados cuantifica numéricamente los efectos del modelo a través de cuatro
escenarios construidos sobre la red de tres firmas (banco nacional $N$, banco local $L$, firma
local $F$) introducida en el capítulo anterior. Los escenarios se diseñan para aislar
progresivamente cada mecanismo:

\begin{enumerate}[label=(\roman*)]
    \item \textbf{3F-WITH:} red única tradicional con la prohibición regulatoria activa ($\Psi$
    desactiva el contrato $N$--$F$).
    \item \textbf{3F-WITHOUT:} red única sin la restricción, que sirve como benchmark de
    bienestar potencial.
    \item \textbf{4F-WITH:} se incorpora la firma Fintech (4 agentes) bajo la misma restricción,
    para aislar el efecto de la competencia entre intermediarios.
    \item \textbf{4F-WITHOUT:} la red de 4 firmas sin restricción, que actualiza el benchmark de
    bienestar potencial con un agente adicional.
\end{enumerate}

La métrica central es la utilidad agregada del sistema, que combina los retornos esperados
netos de pagos y el costo cuadrático del riesgo. Como métricas complementarias se reportan
el volumen bilateral total, las utilidades por agente, las cuotas de mercado y los precios de
equilibrio.

\subsection*{1. Base: Referencia al Ejemplo 1}

El primer experimento compara 3F-WITH y 3F-WITHOUT para medir el costo de la restricción
regulatoria en ausencia de competencia. Cuando se desactiva el contrato $N$--$F$, el volumen
bilateral total cae de $1{,}976$ a $0{,}822$ (una contracción de la actividad de $58{,}4\%$), y la
utilidad agregada del sistema baja de $1{,}434$ a $0{,}706$, lo que equivale a una pérdida de
bienestar de $50{,}8\%$.

La distribución del impacto entre agentes es marcadamente asimétrica. El banco nacional
pierde $84{,}6\%$ de su utilidad (de $0{,}4125$ a $0{,}0636$) y la firma local pierde $66{,}1\%$
(de $0{,}5959$ a $0{,}2020$), porque ambos son los extremos del contrato prohibido. El banco
local, en cambio, es el único que mejora su posición bajo la restricción: su utilidad se mantiene
prácticamente intacta (de $0{,}4254$ a $0{,}4405$, $+3{,}6\%$), porque la prohibición lo
convierte en el único intermediario habilitado entre $N$ y $F$. Las rentas del banco nacional
incluso cambian de signo (de $0{,}0681$ a $-0{,}3598$), reflejando que en lugar de capturar
valor pasa a pagarlo. La cadena de intermediación $N \to L \to F$ transfiere recursos, pero $L$
se queda con el diferencial.

\subsection*{2. Competencia}

El segundo experimento incorpora la firma Fintech (escenario 4F-WITH) manteniendo la
restricción regulatoria. La presencia de un segundo intermediario alternativo al banco local
rompe el monopolio de intermediación: el volumen bilateral sube a $1{,}351$ ($+64{,}3\%$
respecto de la base 3F-WITH) y la utilidad agregada llega a $1{,}034$ ($+46{,}5\%$ sobre la base).

La redistribución entre agentes confirma el mecanismo: el banco nacional gana $159{,}6\%$ de
utilidad y la firma local gana $93{,}6\%$, mientras que el banco local pierde $45{,}7\%$ al perder
su privilegio estructural. Sin embargo, el bienestar agregado de $1{,}034$ sigue $27{,}9\%$ por
debajo del benchmark sin restricción de tres firmas ($1{,}434$). La competencia entre
intermediarios captura parte del bienestar perdido, pero no puede sustituir a la
desregulación. La prohibición sigue activa y el contrato directo $N$--$F$ sigue sin existir.

\subsection*{3. \textit{Dual-network}}

El tercer experimento permite que los agentes asignen contratos en una red alternativa $A$ (con
restricción $\Psi^A$ que sí habilita el contrato $N$--$F$) además de la red tradicional $T$. Antes
de aplicar la prohibición, se verifica el caso simétrico: cuando ambas redes son
regulatoriamente idénticas, el flujo se reparte 50/50 entre $T$ y $A$ y la utilidad agregada
coincide con la del benchmark sin restricción ($1{,}434$ en 3 firmas y $1{,}741$ en 4 firmas).
Esto confirma la Proposición 3.5.7 del modelo: en ausencia de arbitraje regulatorio, el modelo
dual reproduce el modelo base.

Cuando se reintroduce la restricción solo en $T$ (escenario 3F-DUAL-WITH), la red alternativa
absorbe los contratos que $T$ no puede ejecutar. El equilibrio reparte la actividad como
$27{,}6\%$ en $T$ y $72{,}4\%$ en $A$, con el contrato directo $N$--$F$ viajando casi enteramente
por la red alternativa ($W_{13}^A = 1{,}030$). La utilidad agregada del sistema es $1{,}434$,
idéntica al benchmark sin restricción. La red dual recupera el $100\%$ del bienestar perdido por
la prohibición, frente al $27{,}9\%$ de gap que dejaba la sola competencia entre intermediarios.

La distribución por agentes muestra cómo se reparte la ganancia: el banco nacional ve su
utilidad pasar de $0{,}064$ a $0{,}238$ ($+272{,}5\%$), la firma local de $0{,}202$ a $0{,}851$
($+321{,}3\%$), y el banco local cae de $0{,}441$ a $0{,}345$ ($-21{,}7\%$). Es importante notar
que, aun con esta caída, el banco local queda mejor que en el escenario de competencia
4F-WITH ($0{,}239$), porque aquí no compite por flujo dentro de su propia red, sino que conserva
los contratos que $T$ sí permite ejecutar. La diferencia promedio de precios entre redes es
$0{,}059$, lo que puede interpretarse como el \textit{shadow price} de la restricción
regulatoria: cuánto está dispuesto a pagar el sistema por mover un contrato de $T$ a $A$.

El experimento se replica con cuatro firmas (4F-DUAL-WITH): el reparto se vuelve más
equilibrado ($52{,}1\%$ en $T$ y $47{,}9\%$ en $A$) y la utilidad agregada llega a $1{,}741$,
igualando al benchmark ideal de cuatro firmas sin restricción. Frente a la base 3F-WITH la
mejora es de $+146{,}6\%$, y frente al escenario de pura competencia 4F-WITH la mejora es de
$+68{,}4\%$.


% ============================================================
\section*{Oportunidades de Mejora}
% ============================================================

El modelo de Venegas (2024) constituye un avance relevante respecto del marco original de Jalan
y Chakrabarti, principalmente porque introduce de forma tratable la coexistencia de dos
infraestructuras con regulaciones distintas y obtiene un resultado de bienestar limpio. Sin
embargo, esto mismo abre espacio para extensiones futuras.

\subsection*{1. Calibración empírica de los primitivos del modelo}

Las creencias de retorno $\mu_i$, las matrices de covarianza $\Sigma_i$ y los coeficientes de
aversión al riesgo $\gamma_i$ se eligen de manera ilustrativa para sostener los experimentos
numéricos. Esto basta para mostrar mecanismos, pero limita la interpretación cuantitativa de
los resultados: una pérdida de bienestar de $50{,}8\%$ o un gap de $27{,}9\%$ son cifras propias
del modelo, no estimaciones del costo regulatorio en una jurisdicción concreta. Calibrar el
modelo con datos reales permitiría comparar sus resultados con la evidencia empírica, por
ejemplo, contrastando las cuotas $T/A$ que predice con las observadas en jurisdicciones donde
ya coexisten ambas infraestructuras.

\subsection*{2. Incorporar dinámica y dependencia de trayectoria}

En el modelo se caracteriza un equilibrio \textit{pairwise stable} pero no se modela cómo se
llega a él ni cómo evoluciona en el tiempo. El marco teórico discute extensamente situaciones
de \textit{lock-in} y costos de cambio, pero estos elementos no entran al modelo formal. Una
extensión natural sería plantear el problema en tiempo discreto con costos de migración entre
redes, donde la decisión de un agente de moverse a $A$ en $t$ depende de cuántos otros
agentes ya estén en $A$ en $t-1$. Esto permitiría estudiar formalmente en qué condiciones una
red alternativa eficiente en equilibrio puede no ser adoptada por razones de coordinación, y
qué tipo de intervención podría acelerar esta transición.

\subsection*{3. Ausencia de riesgo sistémico y contagio}

El riesgo en el modelo es individual y se captura mediante la covarianza $\Sigma_i$ del portafolio
de cada agente. Eisenberg y Noe (2001) y Elliott, Golub y Jackson (2014), en cambio, ponen el
foco en el contagio entre agentes a través de la red de obligaciones, que es justamente lo que
el modelo no incorpora. Una extensión al presente modelo sería introducir una estructura de
obligaciones mutuas entre los agentes de la red, lo que permitiría comparar la red tradicional
y la alternativa no solo en términos de bienestar agregado sino también de estabilidad
sistémica.

\end{document}